\documentclass[11pt,a4paper]{article}
\usepackage[utf8]{inputenc}
\usepackage[italian]{babel}
\usepackage{amsmath,amssymb}
\usepackage{hyperref}
\usepackage{geometry}
\geometry{margin=2.5cm}
\usepackage{booktabs}
\usepackage{tikz}
\usetikzlibrary{shapes,arrows,positioning}

\title{Il Fantasma: strato di ragionamento interno effimero\\ \large Concept note e prototipo funzionante --- v0.01}
\author{dr.\ Francesco Bulla --- Bologna, 09/10/2026}
\date{Documento locale, bozza non pubblicata --- v0.01}

\begin{document}
\maketitle

\begin{abstract}
Il Fantasma \`e uno strato di ragionamento interno che si attiva prima della generazione
dell'output, produce un pensiero in un codice numerico privato, lo usa per riformulare
la richiesta e poi lo cancella. L'utente vede solo l'output finale. Questa nota raccoglie
l'intuizione originale (Catania, 7 ottobre 2026), tre decisioni di progetto che la rendono
implementabile, il prototipo in \texttt{src/fantasma} e i primi esperimenti live con
DeepSeek-V4.1-Flash (id API \texttt{deepseek-flash}). Quattro casi reali e un test di codice
mostrano che il Fantasma eccelle a incorniciare e riformulare, mentre l'esecuzione resta
compito di un secondo passo. Il trade-off tra privacy cognitiva e verificabilit\`a \`e strutturale
e viene reso esplicito, non nascosto.
\end{abstract}

\section{Idea}
Il Fantasma emerge, agisce, scompare. Lascia traccia solo per chi lo ha evocato.
A differenza di Chain-of-Thought visibile, scratchpad nascosti e latent reasoning,
qui la cancellazione \`e attiva e il codice \`e privato: una sequenza numerica
(es.\ \texttt{0984567345679}) non interpretabile dall'esterno.

\section{Meccanismo}
\begin{enumerate}
\item L'input arriva al modello.
\item Un gate interno decide se attivare il Fantasma.
\item Il Fantasma genera un pensiero numerico privato.
\item Lo stesso modello rilegge il proprio codice (decisione 5.1).
\item L'input riformulato produce l'output finale.
\item Il pensiero viene cancellato; resta solo un impegno crittografico (hash).
\end{enumerate}

\begin{figure}[h]
\centering
\begin{tikzpicture}[node distance=1.1cm, auto,
  blocco/.style={rectangle, draw, rounded corners, text width=2.1cm, align=center, font=\scriptsize},
  freccia/.style={->, thick}]
\node[blocco] (in) {Input};
\node[blocco, right=of in] (gate) {Gate};
\node[blocco, right=of gate] (ghost) {Codice fantasma};
\node[blocco, below=of ghost] (interp) {Stessa AI interpreta};
\node[blocco, left=of interp] (seal) {Sigillo offuscato};
\node[blocco, left=of seal] (ref) {Riformula};
\node[blocco, below=of ref] (out) {Output};
\node[blocco, right=of out] (amn) {Amnesia};
\draw[freccia] (in) -- (gate);
\draw[freccia] (gate) -- (ghost);
\draw[freccia] (ghost) -- (interp);
\draw[freccia] (interp) -- (seal);
\draw[freccia] (seal) -- (ref);
\draw[freccia] (ref) -- (out);
\draw[freccia] (out) -- (amn);
\end{tikzpicture}
\caption{Ciclo del Fantasma: emerge, agisce, scompare. Solo l'audit a hash resta.}
\end{figure}

\section{Tre decisioni di progetto}
\textbf{5.1 --- La stessa AI interpreta.} Nessun decodificatore esterno: il codice
viene riconsegnato allo stesso modello tramite prompt \texttt{FANTASMA-INTERPRETA}.
Solo quello stato, in quel giro, pu\`o leggerlo. Senza salt e nonce non \`e ricostruibile.

\textbf{5.2 --- Cache offuscata.} Il codice grezzo vive solo in memoria effimera per
il giro di interpretazione. Subito dopo viene sigillato in cache in forma offuscata
(shift delle cifre derivato dal salt, cio\`e dallo stato interno). La cache non contiene
mai il grezzo; con salt sbagliato si ottiene rumore.

\textbf{5.3 --- Amnesia dopo l'output.} Sputato l'output, non serve ricordare.
Memoria effimera e cache vengono svuotate. Nessun database vettoriale persistente.
Resta solo l'audit a hash-chain, che registra impegni e non contenuti.

\section{Prototipo}
\texttt{src/fantasma}: \texttt{gate.py} (gate euristico ispezionabile),
\texttt{codec.py} (HMAC + nonce, unicit\`a di ogni Fantasma),
\texttt{memory.py} (\texttt{EphemeralMemory} + \texttt{ObfuscatedCache}),
\texttt{audit.py} (hash-chain senza contenuti),
\texttt{pipeline.py} (flusso 1--8),
\texttt{llm\_deepseek.py} (adapter DeepSeek-V4.1-Flash, chiave solo da env).
Suite: 12 test verdi in locale e su Docker (\texttt{docker run --rm fantasma:0.1}).

\section{Esperimenti live (09/10/2026)}
Modello: DeepSeek-V4.1-Flash (\texttt{deepseek-flash}), \texttt{reasoning\_effort=high},
\texttt{thinking=enabled}. Tutti i casi: ghost attivo, stessa-AI attiva,
memoria e cache vuote dopo, audit verificato.

\begin{center}
\begin{tabular}{lll}
\toprule
Caso & Riformulazione & Output \\
\midrule
Decisione Milano/Catania & valori e vincoli esplicitati & domanda chiarificata \\
Messaggio al collaboratore & tono fermo-professionale & bozza pronta da inviare \\
Brief web app & fasi, alternative, rischi & domanda operativa \\
Paradosso etico & condizioni procedurali & tesi: consenso e contestabilit\`a \\
Test di codice & 4 bug identificati & analisi + fix in 2 passi \\
\bottomrule
\end{tabular}
\end{center}

Sul codice il Fantasma ha individuato tutti i difetti nascosti (default mutabile,
\texttt{range(len+1)}, divisione per zero, nomi bugiardi), prodotto una specifica
eseguibile e, al secondo passo, 391 righe di implementazione con test. Prima esecuzione
rossa per un assert troppo severo (chiavi extra), poi verde dopo un fix di una riga:
\texttt{Tutti i test sono passati.}

\section{Problemi aperti e paradosso etico}
Restano: fondazione del linguaggio privato (Wittgenstein), cancellazione causale
impossibile in un transformer, apprendimento senza memoria persistente, audit senza
rivelazione. Il principio ``nessuno deve decidere come ragiona qualcun altro'' resta
in tensione con il fatto che gate, codice e cancellazione sono progettati: il gate
\`e perci\`o esplicito e configurabile, mai nascosto. Il trade-off privacy--verificabilit\`a
\`e strutturale.

\begin{figure}[h]
\centering
\begin{tikzpicture}[font=\scriptsize]
\draw[->] (0,0) -- (7,0) node[right] {privacy cognitiva};
\draw[->] (0,0) -- (0,4) node[above] {verificabilit\`a};
\draw[thick] (0.5,3.5) .. controls (3,3) and (4,1) .. (6.5,0.5);
\fill (0.7,3.4) circle (2pt) node[above right] {CoT visibile};
\fill (3.4,2.0) circle (2pt) node[above right] {Fantasma};
\fill (6.2,0.6) circle (2pt) node[above] {black box};
\end{tikzpicture}
\caption{Pi\`u il pensiero \`e privato, meno \`e verificabile. Il Fantasma sta nel mezzo, dichiarato.}
\end{figure}

\section{Visione}
Il Fantasma non punta a un modello pi\`u opaco, ma a un'opacit\`a legittima e procedurale:
consenso, contestabilit\`a, limiti temporali, supervisione indipendente del gate.

\begin{figure}[h]
\centering
\begin{tikzpicture}[font=\scriptsize,
  tappa/.style={rectangle, draw, rounded corners, align=center, text width=3.2cm}]
\node[tappa] (v1) {v0.01 \\ prototipo + live};
\node[tappa, right=1cm of v1] (v2) {v0.1 \\ metrica vs CoT nascosto};
\node[tappa, right=1cm of v2] (v3) {v1.0 \\ reset selettivo + supervisione};
\draw[->, thick] (v1) -- (v2);
\draw[->, thick] (v2) -- (v3);
\end{tikzpicture}
\caption{Roadmap: dal prototipo che lavora alla misura, poi alle garanzie.}
\end{figure}

\section{Stato e prossimi passi}
Stato: concept + prototipo v0.01 funzionante, 12 test verdi, esperimenti live registrati.
Prossimi passi: metrica di miglioramento misurabile contro CoT nascosto, reset selettivo
reale, supervisione indipendente del gate. Altrimenti resta nel quaderno: non tutte le
idee devono diventare paper, alcune sono solo belle. Questa, per ora, lavora.
\end{document}
